\documentclass[12pt,a4paper]{article}
\usepackage{amsmath,amssymb}
\usepackage[T1]{fontenc}
\usepackage[utf8]{inputenc}
\title{Laplace Transform Through Its Kernel: A Seven--Step Explanation}
\author{}
\date{}

\begin{document}
	\maketitle
	
	\section*{Overview}
	This note explains precisely what the word ``kernel'' means in the context of the Laplace transform and how the transform works. It presents a seven--step development centered on the integral kernel
	$$
	K(s,t)=e^{-st},
	$$
	and then contrasts this with the algebraic notion of kernel as a null space. Two worked models (an RC circuit and its mechanical analogue) are included.
	
	\section*{ Integral transforms and the Laplace kernel}
	An integral transform has the form
	$$
	\mathcal{T}\{f\}(s)=\int_{0}^{\infty} f(t)\,K(s,t)\,dt,
	$$
	where the function $$K(s,t)$$ is called the \emph{kernel} of the transform. For the one--sided Laplace transform
	$$
	\mathcal{L}\{f\}(s)=F(s)=\int_{0}^{\infty} f(t)\,e^{-st}\,dt,
	$$
	the kernel is
	$$
	K(s,t)=e^{-st}.
	$$
	This use of the word ``kernel'' means the weighting function inside the integral. It is not the same as the kernel (null space) of a linear map.
	
	\section*{Why the kernel is \( e^{-st} \): multiplicativity and damping}
	The kernel satisfies a multiplicativity law in the time variable:
	$$
	K(s,t_1+t_2)=e^{-s(t_1+t_2)}=e^{-st_1}\,e^{-st_2}=K(s,t_1)\,K(s,t_2).
	$$
	Thus, for each fixed $$s$$, the map $$t\mapsto e^{-st}$$ is a character of the additive semigroup $$([0,\infty),+)$$. This property is what turns convolution in time into multiplication in the transform domain.
	
	The same kernel also supplies exponential damping. If $$|f(t)|\le M e^{\alpha t}$$ for large $$t$$, then the integral for $$\Re(s)>\alpha$$ converges because of the factor $$e^{-\Re(s)\,t}$$.
	
	\section*{Core mechanics from the kernel}
	\subsection*{Convolution theorem}
	For causal convolution
	$$
	(f*g)(t)=\int_{0}^{t} f(\tau)\,g(t-\tau)\,d\tau,
	$$
	use the kernel's multiplicativity and Fubini/Tonelli to obtain
	$$
	\mathcal{L}\{f*g\}(s)=\int_{0}^{\infty}\int_{0}^{t} f(\tau)g(t-\tau)\,e^{-st}\,d\tau\,dt
	=\Big(\int_{0}^{\infty} f(\tau)e^{-s\tau}\,d\tau\Big)\Big(\int_{0}^{\infty} g(\theta)e^{-s\theta}\,d\theta\Big)
	=F(s)\,G(s).
	$$
	Hence the Laplace transform is an algebra homomorphism from convolution to pointwise multiplication.
	
	\subsection*{Differentiation and integration}
	For suitable $$f$$ (piecewise differentiable) one has
	$$
	\mathcal{L}\{f'(t)\}
	=\int_{0}^{\infty} f'(t)e^{-st}\,dt
	=\big[f(t)e^{-st}\big]_{0}^{\infty}+s\int_{0}^{\infty} f(t)e^{-st}\,dt
	=-f(0^{+})+sF(s).
	$$
	With zero initial value $$f(0^{+})=0$$, this is $$\mathcal{L}\{f'\}=sF$$. By induction,
	$$
	\mathcal{L}\{f^{(n)}\}=s^{n}F \quad \text{(zero initial conditions)}.
	$$
	Similarly,
	$$
	\mathcal{L}\Big\{\int_{0}^{t} f(\tau)\,d\tau\Big\}=\frac{1}{s}F(s).
	$$
	
	\subsection*{Shifts}
	Time delay by $$a>0$$:
	$$
	\mathcal{L}\{f(t-a)u(t-a)\}=e^{-as}F(s).
	$$
	Exponential weighting by $$e^{at}$$:
	$$
	\mathcal{L}\{e^{at}f(t)\}=F(s-a).
	$$
	Both identities come directly from the factor $$e^{-st}$$ in the kernel.
	
	\section*{Region of convergence and analyticity}
	If $$f$$ is of exponential order $$\alpha$$, then
	$$
	F(s)=\int_{0}^{\infty} f(t)e^{-st}\,dt
	$$
	converges for $$\Re(s)>\alpha$$. On that open half--plane, $$F$$ is holomorphic (complex differentiable). The boundary line $$\Re(s)=\alpha$$ is determined by the growth of $$f$$ and, for meromorphic $$F$$, is related to the location of poles.
	
	\section*{Inversion and uniqueness}
	Under standard hypotheses, $$f$$ can be recovered from $$F$$ via the Bromwich integral
	$$
	f(t)=\frac{1}{2\pi j}\int_{\gamma-j\infty}^{\gamma+j\infty} F(s)\,e^{st}\,ds,
	\quad \gamma \text{ greater than the abscissa of convergence.}
	$$
	For rational $$F$$, residues at the poles yield explicit inverse formulas. This establishes injectivity on the usual classes of signals: in principle, no information is lost.
	
	\section*{What it means to ``collapse \( t \) into \( s \)''}
	The mapping $$f\mapsto F$$ replaces the time variable $$t$$ by the complex parameter $$s=\sigma+j\omega$$. Algebraically:
	\begin{itemize}
		\item Convolution in $$t$$ becomes multiplication in $$s$$.
		\item Differentiation in $$t$$ becomes multiplication by $$s$$.
	\end{itemize}
	Thus a constant--coefficient ODE
	$$
	a_n y^{(n)}+\cdots+a_1 y'+a_0 y=b_m u^{(m)}+\cdots+b_0 u
	$$
	transforms into the algebraic relation
	$$
	(a_n s^{n}+\cdots+a_0)Y(s)=(b_m s^{m}+\cdots+b_0)U(s).
	$$
	The transfer function is the ratio
	$$
	G(s)=\frac{Y(s)}{U(s)},
	$$
	which acts as a multiplier in the transform domain.
	
	\section*{Interpretation via eigenfunctions}
	Exponentials are eigenfunctions of causal LTI systems. If the input is $$e^{st}$$, then the output is $$G(s)e^{st}$$. Decomposing signals against the kernel $$e^{-st}$$ allows one to read off the system's action as multiplication by $$G(s)$$.
	
	\section*{``Kernel'' as null space vs integral kernel}
	In linear algebra, for a linear map $$T:V\to W$$, the kernel (null space) is
	$$
	\ker T=\{v\in V:\, T v=0\}.
	$$
	In canonical decompositions, subspaces such as $$\ker\big((T-\lambda I)^k\big)$$ determine generalized eigenspaces. This notion of kernel is different from the integral kernel.
	
	An integral operator has the form
	$$
	(\mathcal{K}f)(s)=\int K(s,t)f(t)\,dt.
	$$
	Here, $$K$$ is the integral kernel. The null space of the operator is
	$$
	\ker \mathcal{K}=\{f:\, \mathcal{K}f=0\}.
	$$
	For the Laplace operator on standard classes of exponentially bounded signals,
	$$
	\mathcal{L}\{f\}(s)=0 \ \text{for all } s \text{ in a right half--plane} \ \Longrightarrow \ f=0,
	$$
	so the operator's null space is trivial (the transform is injective). This injectivity does not change the meaning of the word ``kernel'' in the integral sense.
	

	
\end{document}
